\chapter{MÔ PHỎNG VÀ ĐÁNH GIÁ}
\label{chap:results}

\section{Mục tiêu và nguyên tắc đánh giá}
Chương này đánh giá sáu phương pháp gồm TD3, DDPG, PPO, AO--SCA, AO--Grid và AnalyticalRIS. Mục tiêu không phải chỉ tìm phương pháp có tổng tốc độ lớn nhất, mà xác định mức cân bằng giữa bốn nhóm chỉ số:
\begin{enumerate}[label=\arabic*)]
  \item tổng tốc độ và xu hướng theo số phần tử STAR--RIS;
  \item tỷ lệ người dùng đạt QoS, xác suất toàn bộ người dùng đạt QoS và mức vi phạm;
  \item độ ổn định giữa các seed huấn luyện;
  \item độ trễ ra quyết định trên cùng một nền tảng CPU.
\end{enumerate}

Toàn bộ kết quả được lấy từ bộ dữ liệu thực nghiệm đã được kiểm tra về tính đầy đủ, tính nhất quán và việc dùng chung tập kịch bản. Các thuật toán DRL dùng tám seed tại mỗi $N$, trong khi mọi phương pháp được đánh giá trên cùng 1.000 kịch bản kiểm thử khóa. Các số liệu không được chọn lọc lại sau khi xem kết quả.

\section{Thiết lập mô phỏng}
\begin{table}[H]
\caption{Thiết lập chính của benchmark sáu phương pháp}
\label{tab:simulation-setup}
\begin{tabularx}{\textwidth}{p{5.1cm}Y}
\toprule
\textbf{Thành phần} & \textbf{Giá trị}\tabularnewline
\midrule
Số người dùng & $K=4$; hai người dùng miền phản xạ và hai người dùng miền truyền qua\\
Số phần tử STAR--RIS & $N\in\{16,32,64,96,128\}$\\
Công suất tối đa & $P_{\max}=1$\\
Công suất nhiễu & $\sigma^2=0{,}001$\\
Ngưỡng QoS & $R_{\min}=0{,}5$ bit/s/Hz\\
Số seed DRL & 8 seed, từ 0 đến 7, cho mỗi thuật toán và mỗi $N$\\
Ngân sách huấn luyện & 100.000 tương tác môi trường cho mỗi seed và mỗi $N$\\
ScenarioBank & 10.000 train, 1.000 validation, 1.000 test cho mỗi $N$\\
Chọn checkpoint & Dựa trên validation, ưu tiên tính khả thi QoS\\
Kiểm thử & Hành động xác định, không có nhiễu khám phá\\
Đo độ trễ & 100 lần lặp cho mỗi phương pháp và mỗi $N$, CPU một luồng\\
\bottomrule
\end{tabularx}
\tablesource{Tác giả tổng hợp từ cấu hình thực nghiệm của luận văn (2026)}
\end{table}

TD3, DDPG và PPO dùng cùng trạng thái, bộ giải mã hành động, reward QoS, ngân sách tương tác và tập dữ liệu. Tuy nhiên, như đã trình bày ở Chương~\ref{chap:method}, siêu tham số không được tìm kiếm với ngân sách bằng nhau cho cả ba thuật toán. Do đó, phần so sánh DRL phản ánh ba cấu hình triển khai đã đánh giá.

\section{Diễn biến trên tập xác thực}
\ThesisFigure{figures/pdf/chapter4_fig01_training_overview.pdf}
{Diễn biến các chỉ số trên tập xác thực trong quá trình huấn luyện TD3, DDPG và PPO}
{fig:ch4-validation-curves}[0.95\linewidth]
\figuresource{Tác giả tổng hợp từ dữ liệu xác thực của ba thuật toán DRL (2026)}

Hình~\ref{fig:ch4-validation-curves} trình bày tổng tốc độ xác thực tại $N=32$ và $N=128$, cùng xác suất toàn bộ người dùng đạt QoS tại $N=128$. Mỗi đường là trung bình tám seed và dải bao quanh biểu diễn một độ lệch chuẩn.

Tại $N=32$, DDPG vẫn cạnh tranh với TD3 và có biến thiên nhỏ hơn ở giai đoạn cuối. Khi tăng lên $N=128$, dải biến thiên của DDPG mở rộng rõ rệt và một số seed suy giảm mạnh. PPO giữ tổng tốc độ quanh mức thấp và không tiến vào vùng toàn bộ người dùng đạt QoS cao trong ngân sách 100.000 tương tác. TD3 duy trì xu hướng ổn định hơn ở cả hai kích thước.

Các đường xác thực không được dùng để kết luận trực tiếp về hiệu năng test; vai trò của chúng là giải thích quá trình lựa chọn checkpoint và sự khác biệt về ổn định. Mọi kết luận định lượng dưới đây dùng tập kiểm thử khóa.

\section{Hiệu năng trên tập kiểm thử}
% Generated from TABLE_SIX_METHOD_PERFORMANCE.csv. Do not edit manually.
\begin{landscape}
\begin{longtable}{r l r r r r r}
\caption{Hiệu năng của sáu phương pháp trên tập kiểm thử khóa.}
\label{tab:six-method-performance}\\
\toprule
$N$ & \textbf{Phương pháp} & \textbf{Số seed} & \textbf{Tổng tốc độ} & \textbf{Tỷ lệ QoS} & \textbf{Toàn bộ UE đạt QoS} & \textbf{Mức vi phạm} \\
\midrule
\endfirsthead
\multicolumn{7}{c}{\tablename\ \thetable\ -- tiếp theo}\\
\toprule
$N$ & \textbf{Phương pháp} & \textbf{Số seed} & \textbf{Tổng tốc độ} & \textbf{Tỷ lệ QoS} & \textbf{Toàn bộ UE đạt QoS} & \textbf{Mức vi phạm} \\
\midrule
\endhead
\bottomrule
\endlastfoot
16 & TD3 & 8 & 10,7661 & 0,9892 & 0,9798 & 0,0082 \\
16 & DDPG & 8 & 8,0956 & 0,9089 & 0,7953 & 0,0436 \\
16 & PPO & 8 & 2,1216 & 0,6242 & 0,0515 & 0,0795 \\
16 & AO--SCA & 1 & 15,2137 & 0,9950 & 0,9830 & 0,0000 \\
16 & AO--Grid & 1 & 3,9814 & 1,0000 & 1,0000 & 0,0000 \\
16 & AnalyticalRIS & 1 & 1,9812 & 0,0000 & 0,0000 & 0,0188 \\
\addlinespace
32 & TD3 & 8 & 10,7160 & 0,9962 & 0,9944 & 0,0025 \\
32 & DDPG & 8 & 10,2295 & 0,9985 & 0,9964 & 0,0020 \\
32 & PPO & 8 & 2,1438 & 0,6468 & 0,0691 & 0,0753 \\
32 & AO--SCA & 1 & 16,4524 & 0,9998 & 0,9990 & $1,27\times 10^{-7}$ \\
32 & AO--Grid & 1 & 3,9819 & 1,0000 & 1,0000 & 0,0000 \\
32 & AnalyticalRIS & 1 & 1,9820 & 0,0000 & 0,0000 & 0,0180 \\
\addlinespace
64 & TD3 & 8 & 11,7024 & 0,9970 & 0,9944 & 0,0027 \\
64 & DDPG & 8 & 9,2887 & 0,8675 & 0,6234 & 0,1564 \\
64 & PPO & 8 & 2,1581 & 0,6638 & 0,0944 & 0,0618 \\
64 & AO--SCA & 1 & 17,8059 & 1,0000 & 1,0000 & 0,0000 \\
64 & AO--Grid & 1 & 3,9820 & 1,0000 & 1,0000 & 0,0000 \\
64 & AnalyticalRIS & 1 & 1,9821 & 0,0000 & 0,0000 & 0,0179 \\
\addlinespace
96 & TD3 & 8 & 12,0636 & 0,9994 & 0,9985 & 0,0005 \\
96 & DDPG & 8 & 7,5368 & 0,6098 & 0,1376 & 0,3972 \\
96 & PPO & 8 & 2,0991 & 0,6102 & 0,0597 & 0,0781 \\
96 & AO--SCA & 1 & 18,8800 & 1,0000 & 1,0000 & 0,0000 \\
96 & AO--Grid & 1 & 3,9821 & 1,0000 & 1,0000 & 0,0000 \\
96 & AnalyticalRIS & 1 & 1,9821 & 0,0000 & 0,0000 & 0,0179 \\
\addlinespace
128 & TD3 & 8 & 12,4117 & 0,9991 & 0,9984 & 0,0010 \\
128 & DDPG & 8 & 4,9659 & 0,5592 & 0,0040 & 0,5512 \\
128 & PPO & 8 & 2,1293 & 0,6138 & 0,0401 & 0,0920 \\
128 & AO--SCA & 1 & 19,6363 & 1,0000 & 1,0000 & 0,0000 \\
128 & AO--Grid & 1 & 3,9821 & 1,0000 & 1,0000 & 0,0000 \\
128 & AnalyticalRIS & 1 & 1,9821 & 0,0000 & 0,0000 & 0,0179 \\
\end{longtable}
\end{landscape}

\noindent\textit{Cách đọc bảng:} Đối với TD3, DDPG và PPO, giá trị được tổng hợp theo tám seed, mỗi seed đánh giá trên 1.000 kịch bản kiểm thử khóa. Đối với AO--SCA, AO--Grid và AnalyticalRIS, độ biến thiên chủ yếu đến từ các kịch bản kênh. Vì hai nhóm có đơn vị bất định khác nhau, không so sánh trực tiếp độ rộng khoảng tin cậy như thể chúng có cùng nguồn biến thiên.


% Vietnamese vector figure generated from results/six_method_v1/tables/TABLE_SIX_METHOD_PERFORMANCE.csv
\begin{landscape}
\begin{figure}[p]
\centering
\begin{tikzpicture}
\begin{groupplot}[
  group style={group size=2 by 2,horizontal sep=1.2cm,vertical sep=0.8cm},
  width=0.41\linewidth,
  height=0.24\linewidth,
  xlabel={Số phần tử STAR--RIS, $N$},
  xtick={16,32,64,96,128},
  grid=both,
  major grid style={line width=.2pt,draw=gray!50},
  minor grid style={line width=.1pt,draw=gray!20},
  tick label style={font=\scriptsize},
  label style={font=\small},
  title style={font=\small\bfseries},
  legend style={font=\scriptsize,draw=none,fill=none}
]

\nextgroupplot[
  title={Tổng tốc độ trung bình},
  ylabel={bit/s/Hz},
  ymin=0,ymax=21,
  legend to name=sixmethodlegend,
  legend columns=6
]
\addplot[blue,mark=*] coordinates {(16,10.7661)(32,10.7160)(64,11.7024)(96,12.0636)(128,12.4117)};\addlegendentry{TD3}
\addplot[teal,mark=square*] coordinates {(16,8.0956)(32,10.2295)(64,9.2887)(96,7.5368)(128,4.9659)};\addlegendentry{DDPG}
\addplot[orange,mark=triangle*] coordinates {(16,2.1216)(32,2.1438)(64,2.1581)(96,2.0991)(128,2.1293)};\addlegendentry{PPO}
\addplot[red,mark=diamond*] coordinates {(16,15.2137)(32,16.4524)(64,17.8059)(96,18.8800)(128,19.6363)};\addlegendentry{AO--SCA}
\addplot[violet,mark=star] coordinates {(16,3.9814)(32,3.9819)(64,3.9820)(96,3.9821)(128,3.9821)};\addlegendentry{AO--Grid}
\addplot[black,mark=x] coordinates {(16,1.9812)(32,1.9820)(64,1.9821)(96,1.9821)(128,1.9821)};\addlegendentry{AnalyticalRIS}

\nextgroupplot[
  title={Tỷ lệ người dùng đạt QoS},
  ylabel={Tỷ lệ},
  ymin=0,ymax=1.05
]
\addplot[blue,mark=*] coordinates {(16,0.9893)(32,0.9962)(64,0.9970)(96,0.9994)(128,0.9991)};
\addplot[teal,mark=square*] coordinates {(16,0.9089)(32,0.9985)(64,0.8675)(96,0.6098)(128,0.5592)};
\addplot[orange,mark=triangle*] coordinates {(16,0.6242)(32,0.6468)(64,0.6638)(96,0.6103)(128,0.6138)};
\addplot[red,mark=diamond*] coordinates {(16,0.9950)(32,0.9998)(64,1.0000)(96,1.0000)(128,1.0000)};
\addplot[violet,mark=star] coordinates {(16,1.0000)(32,1.0000)(64,1.0000)(96,1.0000)(128,1.0000)};
\addplot[black,mark=x] coordinates {(16,0)(32,0)(64,0)(96,0)(128,0)};

\nextgroupplot[
  title={Xác suất toàn bộ người dùng đạt QoS},
  ylabel={Xác suất},
  ymin=0,ymax=1.05
]
\addplot[blue,mark=*] coordinates {(16,0.9798)(32,0.9944)(64,0.9944)(96,0.9985)(128,0.9984)};
\addplot[teal,mark=square*] coordinates {(16,0.7953)(32,0.9964)(64,0.6234)(96,0.1376)(128,0.0040)};
\addplot[orange,mark=triangle*] coordinates {(16,0.0515)(32,0.0691)(64,0.0944)(96,0.0598)(128,0.0401)};
\addplot[red,mark=diamond*] coordinates {(16,0.9830)(32,0.9990)(64,1.0000)(96,1.0000)(128,1.0000)};
\addplot[violet,mark=star] coordinates {(16,1.0000)(32,1.0000)(64,1.0000)(96,1.0000)(128,1.0000)};
\addplot[black,mark=x] coordinates {(16,0)(32,0)(64,0)(96,0)(128,0)};

\nextgroupplot[
  title={Mức vi phạm QoS trung bình},
  ylabel={Mức vi phạm},
  ymin=0,ymax=0.60
]
\addplot[blue,mark=*] coordinates {(16,0.00823)(32,0.00249)(64,0.00269)(96,0.00053)(128,0.00100)};
\addplot[teal,mark=square*] coordinates {(16,0.04357)(32,0.00204)(64,0.15639)(96,0.39717)(128,0.55116)};
\addplot[orange,mark=triangle*] coordinates {(16,0.07947)(32,0.07532)(64,0.06182)(96,0.07807)(128,0.09205)};
\addplot[red,mark=diamond*] coordinates {(16,0.000013)(32,0.0000001)(64,0)(96,0)(128,0)};
\addplot[violet,mark=star] coordinates {(16,0)(32,0)(64,0)(96,0)(128,0)};
\addplot[black,mark=x] coordinates {(16,0.01883)(32,0.01802)(64,0.01794)(96,0.01793)(128,0.01793)};

\end{groupplot}
\end{tikzpicture}
\par\smallskip
\pgfplotslegendfromname{sixmethodlegend}
\caption{So sánh hiệu năng của sáu phương pháp trên cùng tập kiểm thử khóa}
\label{fig:ch4-six-method-performance}
\figuresource{Tác giả vẽ lại từ \texttt{TABLE\_SIX\_METHOD\_PERFORMANCE.csv}, bundle \texttt{six\_method\_v1} đã kiểm toán (2026)}
\end{figure}
\end{landscape}


Hình~\ref{fig:ch4-six-method-performance} được vẽ lại trực tiếp từ bảng hiệu năng đã được kiểm tra, với toàn bộ nhãn bằng tiếng Việt. Bốn panel cho thấy một phương pháp có thể tốt ở một tiêu chí nhưng không tốt ở tiêu chí khác. AO--SCA có tổng tốc độ cao nhất; AO--Grid có QoS hoàn hảo nhưng tổng tốc độ thấp; AnalyticalRIS nhanh và đơn giản nhưng không đạt QoS; TD3 duy trì QoS gần một với tổng tốc độ đứng thứ hai trong phần lớn cấu hình.

\subsection{Kết quả của TD3}
Tổng tốc độ trung bình của TD3 tăng từ 10,7661 bit/s/Hz tại $N=16$ lên 12,4117 bit/s/Hz tại $N=128$. Mức tăng tuyệt đối là 1,6456 bit/s/Hz. Xu hướng này cho thấy Actor khai thác được số phần tử tăng để cải thiện kênh hiệu dụng trong mô hình đã sử dụng.

Tỷ lệ người dùng đạt QoS của TD3 lần lượt là 0,9892; 0,9962; 0,9970; 0,9994 và 0,9991 khi $N$ tăng từ 16 đến 128. Xác suất toàn bộ người dùng đạt QoS nằm từ 0,9798 đến 0,9985. Mức vi phạm trung bình thấp hơn 0,0083 ở mọi cấu hình và giảm xuống khoảng $5,3\times10^{-4}$ tại $N=96$.

Không nên diễn giải xu hướng này thành định luật tăng trưởng điện từ. Bộ sinh kênh không mô hình diện tích bề mặt, path loss theo khoảng cách hoặc aperture. Kết luận phù hợp là TD3 có khả năng mở rộng đến $N=128$ trong môi trường Gaussian và tham số hóa đã sử dụng \cite{Bjornson2020RISMyths,DiRenzo2020SmartRadio}.

\subsection{TD3 so với DDPG}
DDPG đạt tổng tốc độ 8,0956 tại $N=16$, tăng lên 10,2295 tại $N=32$, sau đó giảm còn 9,2887; 7,5368 và 4,9659 tại $N=64,96,128$. Độ lệch chuẩn tổng tốc độ giữa seed tăng lên 5,0875 tại $N=128$, cho thấy độ nhạy lớn đối với quá trình huấn luyện.

Tại $N=32$, DDPG đạt tỷ lệ QoS 0,9985 và xác suất toàn bộ người dùng đạt QoS 0,9964, nhỉnh hơn TD3 tương ứng 0,9962 và 0,9944. DDPG cũng có độ trễ thấp hơn TD3 một lượng nhỏ. Vì vậy, không được kết luận TD3 thống trị DDPG trên mọi chỉ số.

Từ $N=64$ trở lên, cấu hình DDPG được đánh giá suy giảm rõ: xác suất toàn bộ người dùng đạt QoS giảm từ 0,6234 xuống 0,1376 và 0,0040. Tại $N=128$, mức vi phạm trung bình đạt 0,5512. Trong khi đó TD3 giữ xác suất toàn bộ người dùng đạt QoS 0,9984. Trong phạm vi $N\le128$, kết quả quan sát cho thấy cấu hình TD3 ổn định hơn khi kích thước bài toán tăng. Đây là mối liên hệ thực nghiệm; do TD3 và DDPG còn khác về Critic, loss, Layer Normalization, gradient clipping, learning rate và lịch cập nhật, không thể quy sự suy giảm của DDPG chỉ cho số chiều hành động.

\subsection{TD3 so với PPO}
PPO đạt tổng tốc độ khoảng 2,10--2,16 bit/s/Hz trên toàn bộ năm kích thước. Tỷ lệ người dùng đạt QoS nằm khoảng 0,61--0,66 và xác suất toàn bộ người dùng đạt QoS chỉ khoảng 0,04--0,09. PPO vì vậy chưa học được chính sách khả thi trong ngân sách 100.000 tương tác.

Một nguyên nhân hợp lý là PPO là phương pháp on-policy và phải học trên hành động có số chiều từ 57 đến 393. Tuy nhiên, kết quả không chứng minh PPO về bản chất không phù hợp với STAR--RIS--RSMA. Nghiên cứu trước đã cho thấy PPO có thể hoạt động trong một mô hình khác \cite{Meng2024PPOSTARRSMA}. Cần ngân sách tương tác, kiến trúc và tuning riêng để đánh giá giới hạn thực sự của PPO.

\subsection{TD3 so với AO--SCA}
AO--SCA đạt tổng tốc độ 15,2137 tại $N=16$ và tăng lên 19,6363 tại $N=128$, cao nhất trong sáu phương pháp. Chênh lệch AO--SCA so với TD3 tăng từ 4,4476 lên 7,2246 bit/s/Hz. AO--SCA cũng duy trì QoS rất cao: tại $N\ge64$, tỷ lệ QoS và xác suất toàn bộ người dùng đạt QoS đều bằng một trong dữ liệu tổng hợp.

Kết quả này phải được báo cáo nổi bật vì AO--SCA là baseline mạnh nhất về chất lượng. TD3 không thay thế AO--SCA trong bài toán tối ưu offline không giới hạn thời gian. Lợi thế của TD3 nằm ở thời gian ra quyết định, còn khoảng cách chất lượng phản ánh \emph{amortization gap}: Actor tạo nghiệm trong một forward pass thay vì tối ưu trực tiếp từng kịch bản.

AO--SCA vẫn là solver cục bộ phụ thuộc khởi tạo, dùng sai phân hữu hạn và số vòng lặp hữu hạn. Kết quả không được gọi là nghiệm tối ưu toàn cục hoặc cận trên lý thuyết.

\subsection{TD3 so với AO--Grid}
AO--Grid đạt tổng tốc độ gần 3,982 bit/s/Hz ở mọi $N$ và có QoS bằng một. Hành vi này cho thấy codebook tìm được nghiệm bảo thủ, phân phối tài nguyên đủ để thỏa ngưỡng nhưng không khai thác tốt số phần tử tăng để nâng tổng tốc độ.

TD3 cao hơn AO--Grid 6,7848 bit/s/Hz tại $N=16$ và 8,4296 bit/s/Hz tại $N=128$. Dù vậy, AO--Grid minh họa rằng QoS hoàn hảo không đồng nghĩa chất lượng phổ cao. Một codebook dày hơn có thể cải thiện tổng tốc độ nhưng làm tăng số lần đánh giá.

\subsection{TD3 so với AnalyticalRIS}
AnalyticalRIS có tổng tốc độ xấp xỉ 1,982 bit/s/Hz và tỷ lệ QoS bằng không ở mọi $N$. Phương pháp đặt beta, công suất và tốc độ chung bằng nhau, nên không thích nghi với sự khác biệt kênh giữa người dùng. Mức vi phạm khoảng 0,018 và gần như không đổi khi $N$ tăng.

TD3 vượt AnalyticalRIS hơn 8,78 bit/s/Hz tại $N=16$ và hơn 10,43 bit/s/Hz tại $N=128$. Chênh lệch lớn không nên được dùng một mình để khẳng định ưu thế phổ quát, vì AnalyticalRIS là baseline đơn giản nhằm tạo điểm tham chiếu độ trễ.

\section{Đánh giá độ trễ}
% Generated from TABLE_SIX_METHOD_CPU_LATENCY.csv. Do not edit manually.
\begin{landscape}
\begin{longtable}{r l r r r r r}
\caption{Độ trễ ra quyết định của sáu phương pháp trên cùng một CPU một luồng.}
\label{tab:six-method-latency}\\
\toprule
$N$ & \textbf{Phương pháp} & \textbf{Số mẫu} & \textbf{Trung bình (ms)} & \textbf{Độ lệch chuẩn} & \textbf{Trung vị (ms)} & \textbf{Khoảng min--max (ms)} \\
\midrule
\endfirsthead
\multicolumn{7}{c}{\tablename\ \thetable\ -- tiếp theo}\\
\toprule
$N$ & \textbf{Phương pháp} & \textbf{Số mẫu} & \textbf{Trung bình (ms)} & \textbf{Độ lệch chuẩn} & \textbf{Trung vị (ms)} & \textbf{Khoảng min--max (ms)} \\
\midrule
\endhead
\bottomrule
\endlastfoot
16 & TD3 & 100 & 0,246 & 0,011 & 0,243 & 0,233--0,279 \\
16 & DDPG & 100 & 0,220 & 0,012 & 0,215 & 0,203--0,276 \\
16 & PPO & 100 & 0,382 & 0,012 & 0,378 & 0,364--0,407 \\
16 & AO--SCA & 100 & 216,441 & 72,971 & 219,588 & 47,088--319,289 \\
16 & AO--Grid & 100 & 70,603 & 0,384 & 70,576 & 69,988--73,047 \\
16 & AnalyticalRIS & 100 & 0,122 & 0,007 & 0,120 & 0,110--0,151 \\
\addlinespace
32 & TD3 & 100 & 0,258 & 0,011 & 0,255 & 0,223--0,305 \\
32 & DDPG & 100 & 0,232 & 0,011 & 0,230 & 0,203--0,273 \\
32 & PPO & 100 & 0,399 & 0,013 & 0,397 & 0,364--0,435 \\
32 & AO--SCA & 100 & 370,241 & 141,160 & 393,324 & 77,025--547,480 \\
32 & AO--Grid & 100 & 131,349 & 0,428 & 131,306 & 130,523--134,039 \\
32 & AnalyticalRIS & 100 & 0,125 & 0,007 & 0,123 & 0,116--0,161 \\
\addlinespace
64 & TD3 & 100 & 0,285 & 0,011 & 0,282 & 0,266--0,316 \\
64 & DDPG & 100 & 0,258 & 0,011 & 0,254 & 0,240--0,291 \\
64 & PPO & 100 & 0,438 & 0,014 & 0,435 & 0,393--0,476 \\
64 & AO--SCA & 100 & 603,962 & 286,428 & 604,773 & 100,351--1012,014 \\
64 & AO--Grid & 100 & 257,674 & 0,796 & 257,742 & 255,054--259,630 \\
64 & AnalyticalRIS & 100 & 0,135 & 0,009 & 0,133 & 0,112--0,166 \\
\addlinespace
96 & TD3 & 100 & 0,311 & 0,014 & 0,310 & 0,255--0,346 \\
96 & DDPG & 100 & 0,284 & 0,013 & 0,281 & 0,229--0,326 \\
96 & PPO & 100 & 0,487 & 0,014 & 0,486 & 0,428--0,526 \\
96 & AO--SCA & 100 & 890,132 & 441,820 & 905,434 & 151,912--1529,309 \\
96 & AO--Grid & 100 & 399,173 & 1,046 & 399,266 & 396,839--401,746 \\
96 & AnalyticalRIS & 100 & 0,141 & 0,011 & 0,138 & 0,118--0,188 \\
\addlinespace
128 & TD3 & 100 & 0,342 & 0,048 & 0,336 & 0,265--0,777 \\
128 & DDPG & 100 & 0,307 & 0,020 & 0,303 & 0,249--0,396 \\
128 & PPO & 100 & 0,521 & 0,020 & 0,519 & 0,455--0,606 \\
128 & AO--SCA & 100 & 1039,038 & 555,868 & 922,085 & 209,604--2062,616 \\
128 & AO--Grid & 100 & 531,254 & 2,148 & 531,610 & 525,662--535,728 \\
128 & AnalyticalRIS & 100 & 0,150 & 0,009 & 0,148 & 0,123--0,184 \\
\end{longtable}
\end{landscape}

\noindent\textit{Lưu ý:} Các số đo được thực hiện đơn luồng trên cùng một CPU. Đây là độ trễ ra quyết định của thuật toán trong cấu hình đo đã sử dụng, chưa phải độ trễ đầu-cuối của một hệ thống vô tuyến hoàn chỉnh.


\ThesisFigure{figures/pdf/chapter4_fig03_quality_qos_latency_tradeoff.pdf}
{Độ trễ CPU, tỷ lệ so với TD3 và không gian đánh đổi chất lượng--độ trễ}
{fig:ch4-quality-latency}[0.95\linewidth]
\figuresource{Tác giả tổng hợp từ \texttt{TABLE\_SIX\_METHOD\_CPU\_LATENCY.csv} và \texttt{TABLE\_SIX\_METHOD\_PERFORMANCE.csv} (2026)}

\subsection{Độ trễ của nhóm DRL}
Độ trễ trung bình của TD3 tăng từ 0,246 ms tại $N=16$ lên 0,342 ms tại $N=128$. DDPG nhanh hơn TD3 một lượng nhỏ, từ 0,220 đến 0,307 ms; PPO chậm hơn, từ 0,382 đến 0,521 ms, nhưng vẫn ở dưới một mili giây. Phép đo kiểm thử của TD3 và DDPG chủ yếu chạy Actor rồi tính các chỉ số vật lý; hai Critic của TD3 không nằm trên đường suy luận. Vì vậy, chênh lệch nhỏ giữa TD3 và DDPG nên được xem là hệ quả của kiến trúc Actor và overhead triển khai, không phải do TD3 phải chạy hai Critic khi kiểm thử.

Như vậy, TD3 không phải thuật toán DRL nhanh nhất theo số đo này; DDPG có latency thấp hơn. Lợi thế của TD3 là giữ chất lượng và QoS khi $N$ tăng.

\subsection{TD3 so với solver lặp}
AO--Grid có độ trễ từ 70,60 đến 531,25 ms, tương ứng chậm hơn TD3 khoảng 287 đến 1.553 lần. AO--SCA có độ trễ từ 216,44 đến 1.039,04 ms, chậm hơn TD3 khoảng 880; 1.438; 2.121; 2.860 và 3.038 lần tại năm giá trị $N$.

Khoảng cách tăng theo $N$ vì AO--SCA phải ước lượng gradient theo nhiều tọa độ và lặp lại phép tính kênh; AO--Grid phải đánh giá nhiều ứng viên. Actor TD3 tăng gần tuyến tính theo kích thước đầu vào/đầu ra và không thực hiện vòng lặp tối ưu online.

\subsection{AnalyticalRIS và đường biên đánh đổi}
AnalyticalRIS có độ trễ nhỏ nhất, từ 0,122 đến 0,150 ms, nhanh hơn TD3 khoảng hai lần. Tuy nhiên, phương pháp này không đạt QoS trong thí nghiệm. Hình~\ref{fig:ch4-quality-latency} vì vậy thể hiện ba vùng: AnalyticalRIS rất nhanh nhưng chất lượng thấp; AO--SCA chất lượng cao nhưng chậm; TD3 nằm giữa, giữ QoS cao với độ trễ dưới một mili giây.

Các tỷ lệ chỉ áp dụng cho nền tảng CPU một luồng và cấu hình đo đã sử dụng. Chúng không phải cam kết độ trễ đầu-cuối trên mọi phần cứng và chưa bao gồm ước lượng CSI, truyền báo hiệu hoặc thời gian điều khiển phần cứng STAR--RIS.

\section{Phân tích thống kê}
% Generated from results/six_method_v1/tables/TABLE_SIX_METHOD_PAIRED_TESTS_HOLM.csv
% This main-text table keeps comparisons involving TD3 only.

\begin{landscape}
\begin{longtable}{r l r r c c}
  \caption{So sánh ghép cặp tổng tốc độ giữa TD3 và năm phương pháp còn lại sau hiệu chỉnh Holm.}
  \label{tab:td3-paired-tests-holm}\\
  \toprule
  $N$ & \textbf{Phương pháp đối chiếu} & $\Delta=\overline{R}_{\mathrm{TD3}}-\overline{R}_{b}$ & \textbf{Cohen's $d_z$} & \textbf{$t$-Holm} & \textbf{Wilcoxon-Holm} \\
  \midrule
  \endfirsthead
  \multicolumn{6}{c}{\tablename\ \thetable\ -- tiếp theo}\\
  \toprule
  $N$ & \textbf{Phương pháp đối chiếu} & $\Delta=\overline{R}_{\mathrm{TD3}}-\overline{R}_{b}$ & \textbf{Cohen's $d_z$} & \textbf{$t$-Holm} & \textbf{Wilcoxon-Holm} \\
  \midrule
  \endhead
  \bottomrule
  \endlastfoot
  16 & DDPG & 2,6705 & 2,7947 & Có & Có \\
  16 & PPO & 8,6445 & 11,0447 & Có & Có \\
  16 & AO-SCA & -4,4476 & -3,3293 & Có & Có \\
  16 & AO-Grid & 6,7848 & 8,6171 & Có & Có \\
  16 & AnalyticalRIS & 8,7850 & 11,1545 & Có & Có \\
  \addlinespace
  32 & DDPG & 0,4865 & 0,5200 & Có & Có \\
  32 & PPO & 8,5722 & 11,0515 & Có & Có \\
  32 & AO-SCA & -5,7364 & -4,6764 & Có & Có \\
  32 & AO-Grid & 6,7341 & 8,6654 & Có & Có \\
  32 & AnalyticalRIS & 8,7340 & 11,2389 & Có & Có \\
  \addlinespace
  64 & DDPG & 2,4137 & 2,4572 & Có & Có \\
  64 & PPO & 9,5444 & 13,0135 & Có & Có \\
  64 & AO-SCA & -6,1035 & -7,3130 & Có & Có \\
  64 & AO-Grid & 7,7204 & 10,5073 & Có & Có \\
  64 & AnalyticalRIS & 9,7204 & 13,2292 & Có & Có \\
  \addlinespace
  96 & DDPG & 4,5268 & 3,9426 & Có & Có \\
  96 & PPO & 9,9645 & 14,2997 & Có & Có \\
  96 & AO-SCA & -6,8164 & -9,4261 & Có & Có \\
  96 & AO-Grid & 8,0816 & 11,5764 & Có & Có \\
  96 & AnalyticalRIS & 10,0815 & 14,4413 & Có & Có \\
  \addlinespace
  128 & DDPG & 7,4458 & 8,0778 & Có & Có \\
  128 & PPO & 10,2824 & 14,0223 & Có & Có \\
  128 & AO-SCA & -7,2246 & -9,8404 & Có & Có \\
  128 & AO-Grid & 8,4296 & 11,4925 & Có & Có \\
  128 & AnalyticalRIS & 10,4296 & 14,2192 & Có & Có \\
\end{longtable}
\end{landscape}

\noindent Trong bảng, ``Có'' nghĩa là bác bỏ giả thuyết không ở mức ý nghĩa $\alpha=0{,}05$ sau hiệu chỉnh Holm trong cột tương ứng. Dấu của $\Delta$ và $d_z$ phải được đọc cùng nhau: giá trị âm khi so với AO-SCA cho thấy AO-SCA có tổng tốc độ trung bình cao hơn TD3; giá trị dương trong các so sánh còn lại cho thấy TD3 có tổng tốc độ trung bình cao hơn phương pháp đối chiếu. Tại $N=32$, DDPG vẫn có QoS và độ trễ nhỉnh hơn TD3 trên một số chỉ số, vì vậy kết quả tổng tốc độ này không đồng nghĩa TD3 thống trị DDPG trên mọi tiêu chí.

Bảng~\ref{tab:td3-paired-tests-holm} xác định trước một họ 25 giả thuyết liên quan trực tiếp đến TD3, gồm năm phương pháp đối chiếu tại năm giá trị $N$. Đối với kiểm định $t$ ghép cặp, Holm được áp dụng một lần trên toàn bộ 25 $p$-value của họ này; Wilcoxon cũng được hiệu chỉnh độc lập trên đúng 25 $p$-value tương ứng. Sau hiệu chỉnh, cả hai kiểm định đều bác bỏ giả thuyết không ở mức $\alpha=0{,}05$ cho 25/25 trường hợp.

Dấu của chênh lệch phải được đọc cùng kết quả. TD3 có tổng tốc độ cao hơn DDPG, PPO, AO--Grid và AnalyticalRIS tại mọi $N$, nhưng thấp hơn AO--SCA tại mọi $N$. Tại $N=32$, chênh lệch TD3--DDPG chỉ 0,4865 bit/s/Hz và Cohen's $d_z=0,5200$, nhỏ hơn nhiều so với các so sánh khác, dù vẫn có ý nghĩa do 1.000 cặp kịch bản.

Một số giá trị $p$ bị phần mềm lưu thành 0 vì nhỏ hơn giới hạn số học. Trong diễn giải, chúng được xem là nhỏ hơn ngưỡng biểu diễn chứ không phải xác suất bằng không. Ngoài ra, kiểm định dùng kết quả DRL đã trung bình theo seed trên từng kịch bản, nên chưa đồng thời mô hình hóa biến thiên giữa seed và giữa kịch bản.

% Thesis-ready section for Chapter 4.
% Source guardrails: results/six_method_v1/SIX_METHOD_REVIEW.md and SIX_METHOD_AUDIT.json.

\section{Giới hạn của phân tích thống kê và phạm vi diễn giải}
\label{sec:statistical-limitations}

Các kết quả của luận văn được xây dựng trên một giao thức khóa gồm năm kích thước STAR-RIS, tám seed huấn luyện cho mỗi thuật toán học tăng cường sâu và 1.000 kịch bản kiểm thử cho mỗi kích thước. Thiết kế này cho phép so sánh có kiểm soát giữa các phương pháp, nhưng không loại bỏ hoàn toàn bất định do quá trình huấn luyện, mô hình kênh và nền tảng phần cứng. Vì vậy, các kết luận dưới đây chỉ áp dụng cho phạm vi thí nghiệm đã công bố.

\subsection{Đơn vị bất định không hoàn toàn giống nhau giữa hai nhóm phương pháp}

Đối với TD3, DDPG và PPO, các thống kê tổng hợp trước hết được tính trên từng seed rồi mới mô tả sự biến thiên giữa tám seed. Ngược lại, AO-SCA, AO-Grid và AnalyticalRIS là các phương pháp xác định hoặc gần xác định trong giao thức hiện tại, nên độ lệch chuẩn của chúng chủ yếu phản ánh sự khác nhau giữa 1.000 kịch bản kênh. Do đó, khoảng tin cậy của nhóm DRL và nhóm tối ưu truyền thống không đại diện cho cùng một nguồn bất định. Luận văn sử dụng các khoảng này để mô tả kết quả trong từng nhóm, nhưng không diễn giải độ rộng của chúng như một phép so sánh trực tiếp về độ ổn định giữa hai loại phương pháp.

\subsection{Kiểm định ghép cặp chưa phải là mô hình phân cấp đầy đủ}

Các kiểm định ghép cặp sử dụng cùng 1.000 kịch bản kiểm thử. Với các phương pháp DRL, kết quả của tám seed được trung bình theo từng kịch bản trước khi thực hiện kiểm định. Cách làm này tránh xem tám seed như tám bản sao độc lập của cùng một kịch bản, nhưng nó kiểm định sự khác biệt theo kịch bản của một tập hữu hạn các chính sách đã huấn luyện. Phân tích chưa xây dựng mô hình phân cấp đồng thời cho hai nguồn biến thiên là seed huấn luyện và kịch bản kiểm thử. Vì vậy, các giá trị $p$ không được xem là bằng chứng tuyệt đối rằng kết luận sẽ giữ nguyên với mọi seed hoặc mọi phân bố triển khai khác.

\subsection{Hiệu chỉnh nhiều so sánh và hiện tượng tràn số của giá trị $p$}

Bảng kiểm định sử dụng hiệu chỉnh Holm cho họ các so sánh được báo cáo. Một số giá trị $p$ rất nhỏ được phần mềm lưu thành 0 do giới hạn số học dấu phẩy động. Trong luận văn, các trường hợp này phải được viết dưới dạng ``nhỏ hơn ngưỡng biểu diễn'', chẳng hạn $p<10^{-300}$, thay vì khẳng định $p=0$. Bên cạnh giá trị $p$, luận văn đồng thời báo cáo chênh lệch trung bình và Cohen's $d_z$ để thể hiện độ lớn thực tế của hiệu ứng.

\subsection{Hiệu ứng rất lớn có thể xuất phát từ phương sai rất nhỏ}

Một số baseline như AO-Grid và AnalyticalRIS tạo ra kết quả gần như không thay đổi giữa các kịch bản trong thiết lập hiện tại. Khi độ lệch chuẩn của chênh lệch rất nhỏ, Cohen's $d_z$ có thể đạt giá trị rất lớn. Những giá trị này không nên được hiểu theo thang diễn giải thông thường một cách cơ học. Chênh lệch tuyệt đối về tổng tốc độ, QoS và độ trễ vẫn là các đại lượng chính để đánh giá ý nghĩa thực tiễn.

\subsection{Các chỉ số QoS là đại lượng bị chặn}

Tỷ lệ người dùng đạt QoS và xác suất toàn bộ người dùng đạt QoS nằm trong đoạn $[0,1]$. Khoảng tin cậy Student-$t$ không bị chặn có thể vượt nhẹ khỏi miền vật lý khi trung bình gần 0 hoặc 1. Dữ liệu thô được giữ nguyên để bảo đảm khả năng tái lập; khi trình bày trực quan, luận văn giới hạn trục trong $[0,1]$ và không sử dụng phần vượt miền vật lý để đưa ra kết luận.

\subsection{Giới hạn của phép đo độ trễ}

Độ trễ được đo bằng 100 lần lặp cho mỗi phương pháp và mỗi $N$ trên cùng một CPU runner đơn luồng. Thiết kế này giúp so sánh tương đối trong cùng môi trường, nhưng không đại diện cho mọi CPU, GPU, hệ điều hành hoặc thư viện số. Các tỷ lệ tăng tốc chỉ được phát biểu cho nền tảng đã công bố. AnalyticalRIS có độ trễ nhỏ nhất nhưng không đáp ứng QoS; do đó, độ trễ không thể được diễn giải tách rời chất lượng nghiệm.

\subsection{Giới hạn của tính công bằng siêu tham số}

Ba thuật toán TD3, DDPG và PPO dùng cùng ScenarioBank, số tương tác, hàm thưởng QoS, quy tắc chọn checkpoint và tập kiểm thử. Tuy nhiên, thí nghiệm không thực hiện một quá trình tìm kiếm siêu tham số có ngân sách bằng nhau cho mọi thuật toán. TD3 sử dụng Layer Normalization, Huber loss, cắt gradient và chuẩn hóa nhiễu theo chiều action, trong khi DDPG và PPO dùng các thiết lập đơn giản hơn. Vì vậy, kết luận phù hợp là TD3 đạt kết quả tốt nhất trong ba implementation được đánh giá theo giao thức đã công bố; luận văn không khẳng định TD3 luôn vượt mọi cấu hình DDPG hoặc PPO có thể xây dựng.

\subsection{Giới hạn của mô hình và khả năng khái quát}

Các kịch bản kiểm thử được sinh từ mô hình kênh tổng hợp với CSI hoàn hảo, người dùng SISO, vị trí miền truyền và phản xạ cố định, SIC lý tưởng và STAR-RIS thụ động với pha truyền và phản xạ độc lập. Kết quả chưa chứng minh hiệu năng dưới sai số CSI, tương quan không gian thực tế, phần cứng lượng tử hóa, ghép pha vật lý, di động người dùng hoặc dữ liệu đo ngoài hiện trường. AO-SCA cũng chỉ là baseline lặp cục bộ phụ thuộc khởi tạo, không phải nghiệm tối ưu toàn cục hay cận trên lý thuyết.

Tóm lại, kết quả cho phép kết luận rằng TD3 là phương pháp DRL ổn định và có khả năng mở rộng tốt nhất trong ba thuật toán học được đánh giá, đồng thời tạo ra sự cân bằng mạnh giữa tổng tốc độ, độ tin cậy QoS và độ trễ. Kết luận này phải luôn đi kèm phạm vi của giao thức khóa, tám seed huấn luyện, mô hình kênh tổng hợp và nền tảng CPU đã sử dụng.

\section{Khả năng mở rộng}
Khi $N$ tăng từ 16 lên 128, chiều trạng thái tăng từ 172 lên 1.292 và chiều hành động tăng từ 57 lên 393. TD3 vẫn duy trì QoS gần một và tổng tốc độ tăng trong phạm vi khảo sát. Cấu hình DDPG được đánh giá suy giảm khi kích thước bài toán tăng, trong khi PPO giữ hiệu năng thấp; các xu hướng này là quan sát thực nghiệm trên $N\le128$, không phải bằng chứng nhân quả riêng cho số chiều hành động.

MLP hiện tại cần huấn luyện riêng cho từng $N$ và không khai thác tính hoán vị của các phần tử STAR--RIS. Kiến trúc attention, set network hoặc graph network có thể chia sẻ tham số và tổng quát giữa nhiều kích thước. Đây là hướng phát triển cần thiết khi $N$ lớn hơn 128.

\section{Tổng hợp các phát hiện chính}
\subsection{Tính khả thi và QoS của TD3}
TD3 tạo hành động luôn thỏa các ràng buộc hình học nhờ bộ giải mã vật lý. Trên tập kiểm thử khóa, tỷ lệ người dùng đạt QoS từ 0,9892 đến 0,9994 và xác suất toàn bộ người dùng đạt QoS từ 0,9798 đến 0,9985. Kết quả cho thấy cơ chế phần thưởng, điều khiển đối ngẫu và lựa chọn checkpoint duy trì QoS ổn định trong phạm vi mô hình đã xét.

\subsection{So sánh giữa các phương pháp DRL}
TD3 có tổng tốc độ cao hơn DDPG và PPO tại mọi $N$. DDPG cạnh tranh ở $N=32$ và có độ trễ thấp hơn TD3, nhưng mất ổn định khi $N$ tăng. PPO ổn định ở mức hiệu năng thấp nhưng chưa đạt QoS. Vì vậy, TD3 là cấu hình DRL có khả năng mở rộng tốt nhất trong ba phương pháp được đánh giá, không phải bằng chứng rằng TD3 luôn vượt mọi cấu hình DDPG hoặc PPO.

\subsection{So sánh với các phương pháp truyền thống}
AO--SCA đạt tổng tốc độ cao nhất và QoS rất cao nhưng có độ trễ lớn hơn TD3 từ hàng trăm đến hàng nghìn lần tùy $N$. AO--Grid đạt QoS hoàn hảo nhưng tổng tốc độ thấp và cũng chậm hơn TD3 đáng kể. AnalyticalRIS nhanh nhất nhưng không đạt QoS. Kết quả thể hiện rõ đánh đổi giữa chất lượng nghiệm và thời gian ra quyết định.

\subsection{Ảnh hưởng của số phần tử STAR--RIS}
Tổng tốc độ TD3 tăng 1,6456 bit/s/Hz giữa $N=16$ và $N=128$, trong khi QoS được giữ gần một và độ trễ chỉ tăng khoảng 0,096 ms. AO--SCA tăng chất lượng nhiều hơn nhưng độ trễ tăng mạnh; DDPG giảm ổn định; PPO gần như không cải thiện. Xu hướng này chỉ áp dụng cho bộ sinh kênh và cấu hình đã sử dụng.

\section{Hàm ý thực tiễn}
TD3 phù hợp khi hệ thống có thể huấn luyện offline và cần tạo quyết định nhanh cho nhiều hiện thực kênh. Trong tối ưu offline hoặc khi có thời gian giải dài, AO--SCA có thể được ưu tiên. Một hướng triển khai thực tế là dùng Actor TD3 làm điểm khởi tạo nhanh, sau đó chạy một số ít vòng AO--SCA khi ngân sách thời gian cho phép. Cách lai này có thể giảm khoảng cách chất lượng mà vẫn tận dụng ưu thế latency.

Khi phân bố kênh thay đổi, cần phát hiện drift và huấn luyện lại. Khi xét CSI không hoàn hảo hoặc phần cứng ghép pha, trạng thái và bộ giải mã hành động phải được thay đổi thay vì chỉ thêm nhiễu vào mô hình hiện tại.

\section{Kết luận chương}
Benchmark sáu phương pháp cho thấy AO--SCA đạt tổng tốc độ cao nhất, AnalyticalRIS đạt độ trễ thấp nhất, còn TD3 là phương pháp DRL ổn định nhất và duy trì QoS gần một đến $N=128$. TD3 không tốt nhất trên mọi chỉ số, nhưng tạo sự cân bằng rõ giữa chất lượng nghiệm, độ tin cậy QoS và thời gian ra quyết định. Kết luận này được hỗ trợ bởi bảng dữ liệu thô, kiểm định ghép cặp và quy trình kiểm tra tính nhất quán, đồng thời bị giới hạn bởi mô hình SISO, CSI hoàn hảo, kênh tổng hợp và số seed hữu hạn.
